\documentclass[11pt,a4paper]{article}
\usepackage[margin=2cm]{geometry}
\usepackage[T1]{fontenc}
\usepackage{lmodern}
\usepackage{booktabs}
\usepackage{longtable}
\usepackage{array}
\usepackage{enumitem}
\usepackage{xcolor}
\usepackage{amssymb}
\usepackage{xurl}
\usepackage[hidelinks]{hyperref}
\sloppy\emergencystretch=2em
\setlist{itemsep=2pt,topsep=3pt}

\newcommand{\pOne}{\textcolor{red!75!black}{\textbf{P1}}}
\newcommand{\pTwo}{\textcolor{orange!85!black}{\textbf{P2}}}
\newcommand{\pThree}{\textcolor{blue!70!black}{\textbf{P3}}}
\newcommand{\tbd}{\textcolor{red}{\texttt{\textbackslash TBD}}}

\title{CAGI $\rightarrow$ NSS 2026: Remaining Work for the Authors}
\author{Internal guideline, not for submission}
\date{Prepared 27 September 2026 \quad$\cdot$\quad Deadline 30 September 2026, 23:59 AoE (1 October, 18:59 in Vietnam)}

\begin{document}
\maketitle

\section*{0. Status after review, \texttt{/improve} and \texttt{/ready}}
\begin{itemize}[leftmargin=*]
\item \textbf{Reviews.} \texttt{review\_paper\_report.pdf} (Reviewer A, normal mode): Weak Reject, 56/100, confidence 4/5. \texttt{report.tex} (Reviewer B, Q2 standard): Weak Reject, 58/100 as stated, which is 56/100 once B's own 5-point category cap is applied. The two reviews agree on every overlapping point; the consolidated list is in Sect.~11 of Reviewer A's report.
\item \textbf{Done without new data} (details in \texttt{revision\_notes.md}):
  \begin{itemize}
  \item anonymized the paper;
  \item fit it to 15 pages including references;
  \item removed every synthetic or illustrative number;
  \item corrected the ``same learner'' wording and added baseline B3$'$;
  \item framed relative checkpoints as retrospective and moved the post-hoc subgroup out of the abstract and conclusion;
  \item consolidated the limitations;
  \item added and verified 2 references and cited Random Forest;
  \item fixed the internal inconsistencies;
  \item passed every static readiness gate except placeholders.
  \end{itemize}
\item \textbf{What blocks submission:} 19 red \tbd{} boxes in \texttt{cagi\_nss2026.pdf}. They need your result files, which are not in this folder. The macro is now fail-safe. With \verb|\pendingfalse|, any remaining \verb|\TBD| stops the build, so a placeholder cannot slip into the submitted PDF.
\item \textbf{Rule for every task:} copy numbers only from a versioned run (commit hash + config + output CSV). Record the file name next to each value in \texttt{report\_audit.md} (experiment ledger). Report whatever the run gives, even if it weakens a claim; Sect.~3 below says how to word each outcome.
\end{itemize}

\section*{1. Priorities}
\pOne~= blocks submission (placeholder, contradiction, or possible leak). \pTwo~= reviewers will almost certainly ask; do it if time allows before 30 Sep, otherwise for the journal version. \pThree~= after submission.

\section*{2. Tasks}
{\small
\begin{longtable}{@{}>{\raggedright\arraybackslash}p{0.7cm}>{\raggedright\arraybackslash}p{0.6cm}>{\raggedright\arraybackslash}p{6.9cm}>{\raggedright\arraybackslash}p{5.6cm}>{\raggedright\arraybackslash}p{1.25cm}@{}}
\toprule
ID & Pri. & Task, manuscript location, and how to compute & Acceptance check / expected outcome & Due \\
\midrule
\endhead
\multicolumn{5}{@{}l}{\textbf{A. Do first: these can invalidate other numbers}}\\
T0 & \pOne & \textbf{Provenance of the headline.} M1 pooled PR-AUC was 0.736 in the 19 Sep draft (from \texttt{idea.pdf}) and is 0.691 now. Identify the run (commit, config, output directory) that produced 0.691, 0.688, the E4--E7 table and Table~5. Regenerate all tables from that single directory. &
One run ID listed in \texttt{report\_audit.md}; every number in Tables 3--5 and the abstract traced to a file. & 28 Sep \\
T4 & \pOne & \textbf{Prefix-length ratio} (Sect.~4.3 Temporal; Sect.~6.2 case study). Find its formula in the feature code. \emph{Safe} if it uses only $e_1..e_k$ (e.g., $k$ divided by a quantity already observed at step $k$). \emph{Leaks} if its denominator is the final length $n$, the endpoint index, or the relative-checkpoint percentage. &
If safe: write the formula into the \tbd{} in Sect.~4.3. If it leaks: delete the feature, rerun \emph{everything} (Tables 3--5, E4--E7, calibration, Chibi case), and update the abstract numbers. Expected: pooled PR-AUC may fall, and the 25\% checkpoint result is the most likely to change. & 28 Sep \\
T5 & \pOne & \textbf{Chibi case vs.\ Table 5} (Sect.~6.2). The case says the out-of-fold probability is $\approx$0.98 at $k=2$. Table 5 says no incident is alerted at $\theta=0.9$. Check whether 0.98 is raw or Platt-calibrated, whether Chibi is among the 10 endpoint-confirmed incidents, and recompute both from the \emph{same} calibrated out-of-fold file. &
Consistent case and table. Delete the \tbd{} sentence in the case study. If 0.98 is raw, write ``raw score'' and give the calibrated value. & 28 Sep \\
\multicolumn{5}{@{}l}{\textbf{B. Fill the placeholders (no new modelling beyond B3$'$)}}\\
T1 & \pOne & \textbf{RQ1 paired statistics} (Sect.~6.1: \tbd{}\{mean\}, \{CI\}, \{$p$\}). Use the same incident-level bootstrap as Table~4 so the two are comparable: 2{,}000 resamples of the 15 incidents with replacement, seed 42, statistic $=$ PR-AUC(M1) $-$ PR-AUC(B3) on the resampled rows. Two-sided $p=2\min(\Pr[\Delta^*\le0],\Pr[\Delta^*\ge0])$. Report the mean and the 2.5/97.5 percentiles. &
Expected: CI includes 0 (the pooled gap is only +0.003). Keep the current wording ``comparable''. Do \emph{not} use Wilcoxon; Sect.~5.3 now promises a bootstrap $p$. & 29 Sep \\
T3 & \pOne & \textbf{Same-learner baseline B3$'$} (Table~3a row; Sect.~6.1). XGBoost with M1's exact hyperparameters, per-fold \texttt{scale\_pos\_weight}, seed and folds, on the B3 flat feature set. Also run the T1 bootstrap for M1$-$B3$'$; this is the clean RQ1 test. &
Fill the Table 3a cell and the Sect.~6.1 \tbd{}. See Sect.~3 below for wording by outcome. & 29 Sep \\
T2 & \pOne & \textbf{False-alert rates} (Table~5, 3 cells). For each $\theta\in\{0.5,0.7,0.9\}$, count the negative trajectories whose calibrated out-of-fold score reaches $\theta$ at \emph{any} checkpoint. Rate $=$ count$/612$. If the unit is prefix rows instead, say so in the caption. &
Monotone non-increasing in $\theta$. Write as ``$x$\% ($c$/612)''. The previous illustrative values (0.82/0.33/0.33\%) must not be reused unless the run reproduces them. & 29 Sep \\
T6 & \pOne & \textbf{Positive trajectories and prevalence} (Sect.~5.1 \tbd{}; Table~3b, 4 cells). Count positive trajectories (the row count implies $\ge$21). For each relative checkpoint, report the positive rate of the rows it contains. &
Numbers consistent with 3{,}354 rows at 5.0\% ($\approx$168 positive rows). Add one sentence on how multi-trajectory incidents are weighted in the pooled PR-AUC. & 29 Sep \\
T7 & \pOne & \textbf{Spearman with outlier} (Sect.~6.1 \tbd{}). Recompute the per-incident M1$-$B3 gap vs.\ length with all 15 incidents; name the excluded incident and why in the repository. &
If $\rho$ drops sharply, soften ``correlates'' to ``tends to increase'', and keep the paragraph exploratory. & 29 Sep \\
T8 & \pOne & \textbf{Peel-chain thresholds} (Sect.~4.3 Motif, 2 \tbd{}). Copy the forwarding fraction and peel fraction from the config used by the T0 run. &
Two numbers, e.g., ``at least 90\%'' / ``at most 10\%'', matching the code exactly. & 28 Sep \\
T9 & \pOne & \textbf{Chibi attributions} (Sect.~6.2 \tbd{}). After T4, extract the top-3 TreeSHAP features for Chibi at $k=2$ from the final model. &
Replace the \tbd{} with the actual third feature (or all three, if the order changes). & 29 Sep \\
\multicolumn{5}{@{}l}{\textbf{C. Artifact, final gate, submission}}\\
T10 & \pOne & \textbf{Repository contents.} Sect.~7 now promises the incident registry, decoded datasets, splits, \emph{per-incident results}, and scripts that regenerate every table. Upload all of them anonymously; scrub names, e-mails, paths and git author metadata. &
The link opens in a private window, and \texttt{make tables} (or equivalent) reproduces Tables 3--5 bit-for-bit. & 29 Sep \\
T11 & \pOne & \textbf{Final gate.} Replace every \verb|\TBD{...}| by its number; set \verb|\pendingfalse|; run \texttt{latexmk -pdf -halt-on-error}. The build \emph{fails} if any TBD remains. Then run \texttt{pdfinfo} ($\le$15 pages, empty Author); \texttt{pdftotext | grep -i tbd} (empty); rerun \texttt{audit\_manuscript.py --page-limit 15}; view every page. &
0 TBD, 0 undefined refs, 0 overfull boxes, $\le$15 pages, anonymous. Update \texttt{report\_audit.md} to READY. & 30 Sep AM \\
T12 & \pOne & \textbf{Submit} on EasyChair (\url{https://easychair.org/my/conference?conf=nsssocialsec2026}) by 30 Sep 12:00 Vietnam time; download the uploaded PDF and open it. & Submission ID recorded. & 30 Sep \\
\multicolumn{5}{@{}l}{\textbf{D. Strengthening (reviewers will ask)}}\\
T13 & \pTwo & \textbf{Seed robustness}: rerun M1, B3, B3$'$ with 5 seeds and report mean $\pm$ s.d.\ of pooled PR-AUC (one sentence in Sect.~6.1). & Expected s.d.\ $<0.02$; otherwise the +0.003 gap is noise by construction. & 30 Sep / journal \\
T14 & \pTwo & \textbf{Nested LOIO for feature selection}: inside each training split, choose with vs.\ without the bridge-context features by inner LOIO, then score the held-out incident. Removes the optimism now disclosed in Sect.~6.1 and 7. & Report the nested PR-AUC next to 0.691; if the gap is small, delete the optimism caveat. & journal \\
T15 & \pTwo & \textbf{Online horizons}: PR-AUC and alert rate at 1, 5, 15, 30 minutes after the seed event and at fraction of seed value moved (both available live). Replaces the retrospective 25/50/75\% diagnostic as the deployment evidence. & New small table; expected lower than 0.933 but well above the 0.05 floor. & journal \\
T16 & \pTwo & \textbf{Calibration uncertainty}: incident-clustered bootstrap CI for Brier and ECE; a reliability diagram. & CI reported in Sect.~6.3. & journal \\
T17 & \pTwo & \textbf{Negative sampling}: publish tolerances for chain/period/value/activity matching, the second-pass rule for the 107 complex negatives, the reason for the 5 manual controls, and per-incident negative counts. & Deterministic script in the repository. & journal \\
T18 & \pTwo & \textbf{Trajectory-builder validity}: share of verified (nonce/message-hash) vs.\ amount/timing bridge links; precision on a hand-checked sample of $\ge$50 links; sensitivity to hop limit $\{4,6,8\}$ and retained value $\{2.5,5,10\}\%$. & Small table; stable PR-AUC across settings supports Sect.~4.2. & journal \\
T19 & \pThree & \textbf{Controlled representation study}: same learner on Flat, Flat+Types, Flat+Motifs, Flat+Types+Motifs; plus a small GRU/1-D CNN over typed action tokens under LOIO. & Answers ``is it typing, motifs, or the learner?'' (Reviewer B, M8). & journal \\
T20 & \pThree & \textbf{Mimicry stress test} (insert $j$ benign transfers between laundering steps) and \textbf{per-bridge-family} detection breakdown. & Discussion material. & journal \\
T21 & \pThree & \textbf{More incidents} ($N\ge30$), with the long-vs-short hypothesis (split at 40 actions) \emph{pre-registered} before looking at the new data. & Confirms or refutes the exploratory finding. & journal \\
\bottomrule
\end{longtable}
}

\section*{3. How to word RQ1 after T1 and T3}
\begin{center}\small
\begin{tabular}{@{}p{6.2cm}p{9.6cm}@{}}
\toprule
Outcome & Wording in abstract / Sect.~6.1 / conclusion \\
\midrule
M1$-$B3$'$ CI includes 0 (most likely) & Keep the current text: ``comparable to a flat-graph baseline''; the typing effect is a hypothesis. No change to the abstract. \\
M1$-$B3$'$ $>0$, CI excludes 0 & Add to the abstract: ``and exceeds the same learner on untyped features by $x$ (95\% CI $[\cdot,\cdot]$)''. Restore a typing claim in contribution 3. Do \emph{not} add the post-hoc subgroup back. \\
B3$'$ $>$ M1 & State it plainly in Sect.~6.1 and reframe the contribution around the protocol and calibrated alerts (already the current framing). \\
T4 reveals a leak & Rerun first; all three rows above then apply to the new numbers. \\
\bottomrule
\end{tabular}
\end{center}

\section*{4. Placeholder map (19 boxes in the PDF)}
\begin{center}\small
\begin{tabular}{@{}lll@{}}
\toprule
Location & Placeholder & Task \\
\midrule
Sect.~4.3 Temporal & prefix-length-ratio formula & T4 \\
Sect.~4.3 Motif & forwarding fraction; peel fraction (2) & T8 \\
Sect.~5.1 & number of positive trajectories & T6 \\
Sect.~6.1 ¶1 & PR-AUC of B3$'$; mean; CI; $p$ (4) & T3, T1 \\
Sect.~6.1 ¶2 & $\rho$, $p$ with outlier included & T7 \\
Table~3a & B3$'$ PR-AUC & T3 \\
Table~3b & positive rate at 25/50/75/100\% (4) & T6 \\
Table~5 & false-alert rate at $\theta=0.5/0.7/0.9$ (3) & T2 \\
Sect.~6.2 case study & calibrated or raw; Table~5 consistency & T5 \\
Sect.~6.2 case study & third top attribution & T4, T9 \\
\bottomrule
\end{tabular}
\end{center}

\section*{5. Schedule}
\begin{center}\small
\begin{tabular}{@{}lp{13cm}@{}}
\toprule
Date (Vietnam time) & Work \\
\midrule
Sun 28 Sep & T0 provenance; T4 prefix-length ratio (if it leaks, start the full rerun immediately); T8 thresholds; T5 Chibi vs.\ Table~5. \\
Mon 29 Sep & T3 B3$'$; T1 and T3 bootstraps; T2 false-alert rates; T6 counts and prevalence; T7; T9; T10 repository. If time allows, T13 seeds. \\
Tue 30 Sep, AM & T11 final gate (fill TBDs, \verb|\pendingfalse|, build, audit, view every page). Recompile \texttt{cagi\_nss2026\_marked.tex} if you want to keep an internal diff. \\
Tue 30 Sep, by 12:00 & T12 submit. The hard deadline is 1 Oct 18:59 Vietnam time; keep the rest as buffer. \\
After submission & T14--T21 for the extended journal version. \\
\bottomrule
\end{tabular}
\end{center}

\section*{6. Files in this folder}
\begin{itemize}[leftmargin=*]
\item \texttt{cagi\_nss2026.tex/.pdf}: clean submission copy (anonymous, 15 pp., \verb|\pendingtrue|).
\item \texttt{cagi\_nss2026\_marked.tex/.pdf}: changes since the 27 Sep draft in blue.
\item \texttt{cagi\_nss2026\_before\_improve.tex}: the 27 Sep draft as received (named authors; synthetic values).
\item \texttt{review\_paper\_report.tex/.pdf}: Reviewer A plus the consolidated assessment with \texttt{report.tex}.
\item \texttt{revision\_notes.md}: concern-by-concern change log. \texttt{report\_audit.md} and \texttt{report\_audit.static.md}: readiness audit.
\item \texttt{Todo-list\_2026-09-19.tex}: the previous to-do list, kept for reference.
\end{itemize}

\end{document}
